% 자동 생성 — tools/make_paper_tables.py. 손으로 고치지 말 것.
\begin{table}[t]
\centering
\caption{교전 시나리오와 관측 설정. 값은 학습 체크포인트의 설정에서 직접 읽었다.}
\label{tab:scenario}
\footnotesize
\resizebox{\ifdim\width>\linewidth\linewidth\else\width\fi}{!}{%
\begin{tabular}{lrl}
\toprule
항목 & 값 & 단위 \\
\midrule
적 USV 수 $M$ & 10 & 척 \\
방어정 수 $P$ & 3 & 척 \\
적 속력 & 9.0 & m/s \\
방어정 속력 & 6.0 & m/s \\
교전 해역 한 변 & 12600 & m \\
그물벽 최대 길이 & 450 & m \\
관측 격자 반폭 & 6300 & m \\
관측 격자 해상도 & 50 & 칸 \\
적 무리 최대 수 $C$ & 4 & 개 \\
관측 픽셀 한 변 & 252 & m \\
방어정당 그물 (학습 / 평가) & 1 / 3 & 장 \\
정책 파라미터 수 & 244,389 & 개 \\
\bottomrule
\end{tabular}
}
\\[3pt]
\begin{minipage}{\linewidth}\footnotesize
방어정 속력이 적 속력보다 낮다 --- 추격이 성립하지 않으므로 접근 회랑의 선제 차단이 유일한 대응이다. 그물 보유량은 학습 시 1장, 평가·배포 시 3장이다: 정책은 결정마다 그물벽 하나를 놓으므로 결정당 행동공간이 보유량과 무관하고, 1장으로 학습한 정책이 재학습 없이 3장 설정에서 동작한다.
\end{minipage}
\end{table}
